Soundness (Section~\ref{sec:soundness}) says a marked term is a unitary;
it does not say the machine register can be discarded. This section and
Definition~\ref{def:scratch} separate five assertions. A \emph{clean
implementation} restores only its workspace $Z$, without a reset promise
for logical records. A \emph{clean unitary} returns $M$ to its exact
initial vector, $WE=EA$, so $M$ may be dropped and the induced operation
on $Q$ is the unitary $A$. \emph{Reset of a factor $F$ on a subspace
$\mathcal S$} means only that $F$ is $\ket0$ after $W$ for inputs in
$\mathcal S$---the right notion inside a larger computation, where
$\mathcal S$ is whatever the preceding steps actually reach.
\emph{Marginal restoration} means only that $M$'s reduced state is
unchanged, which permits neither deletion nor replacement by fresh zeros
when $M$ is correlated with data. Finally the reduced operation on $Q$
can be unitary while $M$ ends in a fixed vector $\eta\neq m_0$ ((SA2)):
discarding $M$ is then harmless for $Q$ alone, but $M$ has not been reset
and cannot be reused as blank. Theorems~\ref{thm:sA-clean} and
\ref{thm:sA-recursive} give exact tests for logical restoration, subspace
reset and marginal restoration; Counterexample~\ref{ex:sA-rec-limits}
shows the limits of their implications. The collective disposal test for
the dynamic interpreter in Theorem~\ref{thm:tB-clean} has the same
fixed-pure-factor scope as (SA2).

\subsection{Exact tests for cleanliness and for a unitary reduced channel}

\begin{theorem}[Exact cleanliness]\label{thm:sA-clean}
Let $Q,M$ be finite registers, $\ket{m_0}_M$ a fixed unit vector,
$E\psi=\psi\otimes\ket{m_0}$ and $P_0=EE^\dagger$ as in
Definition~\ref{def:clean}. Let $U,K$ be total unitaries on $QM$, and
put $W=U^\dagger KU$, $\Pi=UP_0U^\dagger$. The following are equivalent:
\begin{enumerate}
 \item $WE=EA$ for a unitary $A$ on $Q$;
 \item $(\id-P_0)WE=0$;
 \item $(\id-\Pi)K\Pi=0$;
 \item $[K,\Pi]=0$, where $[K,\Pi]=K\Pi-\Pi K$.
\end{enumerate}
When these hold, $A=E^\dagger WE$ is unique. Equivalently, for every
input density operator $\rho$ on $Q$, the final state
$WE\rho E^\dagger W^\dagger$ has $M$ marginal $\proj{m_0}$.
For a single input density operator $\rho$, this return condition is exactly
\[
 \tr\bigl[(\id-P_0)WE\rho E^\dagger W^\dagger\bigr]=0.       \tag{SA1}
\]
For every finite untouched reference $R$, the reference forms use
$E\otimes\id_R$, $W\otimes\id_R$, $A\otimes\id_R$ and
$P_0\otimes\id_R$, with factors reordered by their names. In (SA1)
the input may then be any density operator on $QR$, including an
entangled one. The universal equivalences are unchanged.

The channel $\Phi(\rho)=\tr_M(WE\rho E^\dagger W^\dagger)$ on $Q$
is unitary if and only if there are a unitary $A$ on $Q$ and a fixed
unit vector $\eta\in\cH_M$ such that
\[
 WE\psi=A\psi\otimes\eta\quad\hbox{for every }\psi\in\cH_Q.
                                                               \tag{SA2}
\]
Here $\eta$ need not be $m_0$: a unitary reduced channel need not be clean.
If (SA2) holds, then tensoring it with $\id_R$ gives
$(W\otimes\id_R)(E\otimes\id_R)\Psi=(A\otimes\id_R)\Psi\otimes\eta$ for
every $\Psi\in\cH_{QR}$, and the reduced channel on $QR$ is conjugation
by $A\otimes\id_R$. Conversely, items 1–4 and the (SA2) equivalence for
the tensored data follow by applying this theorem with $QR$ in place of
$Q$ and $U\otimes\id_R$, $K\otimes\id_R$ in place of $U$, $K$; the
uniqueness in item 1 then forces the induced unitary on $QR$ to be
$A\otimes\id_R$.
\end{theorem}
\begin{proof}
 \step{sA-clean-restrict}{Conditions 1 and 2 are equivalent, with the
 stated unique $A$.}
 \begin{proof}
  \step{sA-clean-12}{\assume{Condition 1.}\prove{Condition 2.}}
  \begin{proof}\pf\ $(\id-P_0)E=0$, so multiply $WE=EA$ by
  $\id-P_0$ on the left.
  \end{proof}
  \step{sA-clean-21}{\assume{Condition 2.}\prove{Condition 1.}}
  \begin{proof}
   \step{sA-clean-A}{Set $A=E^\dagger WE$. Then $EA=WE$ and
   $A^\dagger A=\id_Q$.}
   \begin{proof}\pf\ Condition 2 says $P_0WE=WE$, whence
   $EA=P_0WE=WE$ and
   $A^\dagger A=E^\dagger W^\dagger P_0WE=E^\dagger E=\id_Q$.
   \end{proof}
   \qedstep
   \begin{proof}\pf\ A square isometry on a finite-dimensional space
   is surjective: its kernel is zero, so its rank equals the dimension.
   Thus also $AA^\dagger=\id_Q$. Multiplication by $E^\dagger$ shows
   uniqueness.
   \end{proof}
  \end{proof}
  \qedstep
  \begin{proof}\pf\ By \stepref{sA-clean-12} and
  \stepref{sA-clean-21}.
  \end{proof}
 \end{proof}
 \step{sA-clean-conjugate}{Conditions 2 and 3 are equivalent.}
 \begin{proof}\pf\ The range of $E$ is the range of $P_0$, so
 condition 2 is equivalent to $(\id-P_0)WP_0=0$. Conjugate this
 equality by $U$ and substitute $W=U^\dagger KU$ to obtain
 $(\id-\Pi)K\Pi=0$. Conjugation is invertible.
 \end{proof}
 \step{sA-clean-commute}{Conditions 3 and 4 are equivalent.}
 \begin{proof}
  \step{sA-clean-invariant}{\assume{Condition 3.}
  \prove{$K\operatorname{ran}\Pi=\operatorname{ran}\Pi$.}}
  \begin{proof}\pf\ Condition 3 gives inclusion. The restriction of
  the unitary $K$ is injective, so its image has the same finite
  dimension as $\operatorname{ran}\Pi$, proving equality.
  \end{proof}
  \step{sA-clean-reduce}{Under condition 3 the orthogonal complement
  of $\operatorname{ran}\Pi$ is invariant too.}
  \begin{proof}\pf\ For $x\perp\operatorname{ran}\Pi$ and
  $y\in\operatorname{ran}\Pi$, write $y=Kz$ with
  $z\in\operatorname{ran}\Pi$ using \stepref{sA-clean-invariant}.
  Then $\langle y,Kx\rangle=\langle z,x\rangle=0$.
  Thus $K$ preserves both summands of the orthogonal decomposition,
  and commutes with their projection $\Pi$.
  \end{proof}
  \qedstep
  \begin{proof}\pf\ This proves 3 implies 4. Conversely
  $K\Pi=\Pi K$ implies $(\id-\Pi)K\Pi=0$ by multiplication.
  \end{proof}
 \end{proof}
 \step{sA-clean-one}{Equation (SA1) is equivalent to return of $M$
 for that input.}
 \begin{proof}
  \step{sA-clean-positive}{For a positive operator $\sigma$, and an
  orthogonal projection $P$, $\tr[(\id-P)\sigma]=0$ if and only if
  $\operatorname{supp}\sigma\subseteq\operatorname{ran}P$.}
  \begin{proof}\pf\ A spectral decomposition
  $\sigma=\sum_j\lambda_j\proj{v_j}$ with $\lambda_j>0$ on its
  support gives the trace as
  $\sum_j\lambda_j\|(\id-P)v_j\|^2$. Its nonnegative summands
  vanish precisely when every $v_j$ lies in that range.
  \end{proof}
  \qedstep
  \begin{proof}\pf\ Apply \stepref{sA-clean-positive} to
  $\sigma=WE\rho E^\dagger W^\dagger$ and $P=P_0$.
  Support in $\cH_Q\otimes\ket{m_0}$ is equivalent to marginal
  $\proj{m_0}$: one direction follows by partial trace; the other
  follows from
  $\id-P_0=\id_Q\otimes(\id_M-\proj{m_0})$ and
  $\tr[(\id-P_0)\sigma]
  =\tr[(\id_M-\proj{m_0})\tr_Q\sigma]$.
  Testing all pure inputs yields condition 2. Tensoring these
  identities with $\id_R$ gives the stated reference form; the
  implication back to the unreferenced test uses a one-dimensional $R$. Conversely, items 1--4 and the (SA2) equivalence for the tensored data follow by applying this theorem with $QR$ in place of $Q$ and $U\otimes\id_R$, $K\otimes\id_R$ in place of $U$, $K$; the uniqueness in item 1 then forces the induced unitary on $QR$ to be $A\otimes\id_R$.
  \end{proof}
 \end{proof}
 \step{sA-clean-channel}{The reduced channel is unitary exactly under (SA2).}
 \begin{proof}
  \step{sA-clean-channel-if}{\assume{Equation (SA2).}
  \prove{$\Phi(\rho)=A\rho A^\dagger$.}}
  \begin{proof}\pf\ Expand $\rho$ in matrix units and trace the
  fixed factor $\proj\eta$, whose trace is one.
  \end{proof}
  \step{sA-clean-channel-only}{\assume{$\Phi(\rho)=A\rho A^\dagger$
  for a unitary $A$ and every $\rho$.}\prove{Equation (SA2).}}
  \begin{proof}
   \step{sA-clean-basis}{\pick\ an orthonormal basis
   $e_1,\ldots,e_d$ of $Q$. With $V=(A^\dagger\otimes\id_M)WE$,
   there are unit vectors $\eta_j$ such that $Ve_j=e_j\otimes\eta_j$.}
   \begin{proof}\pf\ A finite Hilbert space has an orthonormal basis.
   The $Q$ marginal of $\proj{Ve_j}$ is $\proj{e_j}$; the support
   argument of \stepref{sA-clean-one}, with the factors exchanged,
   forces $Ve_j$ into $\ket{e_j}\otimes\cH_M$. Norm preservation
   makes $\eta_j$ unit.
   \end{proof}
   \step{sA-clean-same}{Every $\eta_j$ is the same vector.}
   \begin{proof}\pf\ If $i\ne j$, apply the unchanged $Q$ marginal
   property of $V$ to $(e_i+e_j)/\sqrt2$. Its $(i,j)$ matrix element
   is $\langle\eta_j,\eta_i\rangle/2$, and must be $1/2$.
   Hence $\|\eta_i-\eta_j\|^2=2-2\operatorname{Re}
   \langle\eta_j,\eta_i\rangle=0$. For $d=1$ there is just one
   vector and the assertion already holds.
   \end{proof}
   \qedstep
   \begin{proof}\pf\ Set $\eta=\eta_1$ and extend the basis
   equality by linearity; multiply $V$ by $A\otimes\id_M$.
   \end{proof}
  \end{proof}
  \qedstep
  \begin{proof}\pf\ By \stepref{sA-clean-channel-if} and
  \stepref{sA-clean-channel-only}.
  \end{proof}
 \end{proof}
 \qedstep
 \begin{proof}\pf\ The four equivalences follow from
 \stepref{sA-clean-restrict}--\stepref{sA-clean-commute}; the
 input and channel assertions follow from \stepref{sA-clean-one}
 and \stepref{sA-clean-channel}.
 \end{proof}
\end{proof}

\begin{definition}[Logical occurrences and retained segments]
\label{def:sA-reset-points}
For a finite fixed-space term define two annotated ordered expansions
$\mathcal E_+(P)$ and $\mathcal E_-(P)$. Lists here are in execution
order; the minus sign denotes the adjoint. Identity gives an empty list.
For an atom $a$ the lists are $(A_a)$ and $(A_a^\dagger)$.
For a sequence, concatenate in the orders
\[
 \mathcal E_+(R;S)=(\mathcal E_+(R),\mathcal E_+(S)),\qquad
 \mathcal E_-(R;S)=(\mathcal E_-(S),\mathcal E_-(R)).
\]
For each expanded occurrence of a marked term set
\[
 \mathcal E_\epsilon(\mathsf M(R,K))
   =[\mathcal E_+(R),K_\epsilon,\mathcal E_-(R)]_{o,\epsilon},
 \qquad K_+=K,\quad K_-=K^\dagger,
 \quad\epsilon\in\{+,-\}.
\]
The brackets are annotations, not operations: $o$ distinguishes the
individual occurrence in the finite expansion, even when several
occurrences come from the same syntactic subterm. Its entry is immediately
before its bracketed list, its exit immediately after. A \emph{marked
reset point} is that exit together with its annotated occurrence and
a specified retained factor $F_o$ to test for zero. The atom and record
identifiers stay those of the original syntax in both expansions.
Outer allocations are left as zero embeddings, outer sequences are
concatenated, and each run is replaced by its plus expansion.
Atoms, including derived menu atoms, are indivisible at this logical
level: no marked occurrences inside their supplied implementations or
inside the menu entries are enumerated here.

More generally, in a specified retained circuit of total unitary gates
and fresh zero embeddings, a \emph{segment} is a consecutive list of
operations with entry and exit boundaries, possibly the empty list.
Its product is an isometry $V:\cH_A\to\cH_B$ on its typed input
and output lists. This definition also applies to a supplied physical
implementation, using its full register lists including workspace.
Thus a boundary inside an implemented atom or keep has a meaning only
after that implementation has been specified. A reset test certifies
the state already produced there; it inserts no operation and deletes
no register.
\end{definition}

\begin{theorem}[Recursive reset and marginal-restoration tests]
\label{thm:sA-recursive}
Let $L$ be a finite retained register list, $U,K$ total unitaries on it,
$W=U^\dagger KU$, and $\mathcal S\subseteq\cH_L$ a specified subspace
with orthogonal projection $P_{\mathcal S}$. Let $F\subseteq L$ be the
auxiliary factor proposed for reset, and set
$P_F=\id_{L\setminus F}\otimes\proj0_F$ (the projector $P_A$ of
Definition~\ref{def:clean}, with $F$ in place of $A$). For arbitrary sublists,
these tensor expressions use the canonical permutation of factors from
their order in $L$ into the displayed order, and are transported back
by its inverse. The same convention applies to $M'$ below.
Then $F$ is zero after $W$ on every input
supported on $\mathcal S$ if and only if
\[
 (\id-UP_FU^\dagger)KU P_{\mathcal S}=0.                  \tag{SA3}
\]
For reuse as an initially blank record the additional incoming condition
$\mathcal S\subseteq\operatorname{ran}P_F$ is required. If
$\mathcal S=\operatorname{ran}P_F$, (SA3) is equivalent to
$[K,UP_FU^\dagger]=0$.

Let $M'\subseteq L$ instead be a factor whose marginal is to be restored,
with no product-state assumption. Fix a basis of $M'$ and its matrix
units $B_{ij}=\ket i\bra j$. Preservation of the joint reduced state
on $M'$ and every finite untouched external reference $R$, for every
joint density operator supported on $\mathcal S\otimes\cH_R$ under
$W\otimes\id_R$, is equivalent to
\[
 P_{\mathcal S}W^\dagger(\id_{L\setminus M'}\otimes B_{ij})
 WP_{\mathcal S}
 =P_{\mathcal S}(\id_{L\setminus M'}\otimes B_{ij})
 P_{\mathcal S}\quad\hbox{for all }i,j.                 \tag{SA4}
\]

For any specified segment isometry $V:\cH_A\to\cH_B$, input
subspace $\mathcal S\subseteq\cH_A$, and factor $F\subseteq B$,
set $P_F^B=\id_{B\setminus F}\otimes\proj0_F$. Reset at its exit
on every input supported on $\mathcal S$ is equivalent to
\[
 (\id_B-P_F^B)VP_{\mathcal S}=0.                       \tag{SA3a}
\]
This includes physical segment boundaries inside atoms, with all
workspace included in $A,B$ and no zero-workspace assumption there.

For a finite Tier~A term or typed outer program, also with the derived
menus of Definition~\ref{def:tierb}, fix an initial subspace
$\mathcal S_0$ and a finite collection of marked reset points as
defined in Definition~\ref{def:sA-reset-points}. Marked occurrences inside the entries of a derived menu atom are not
marked reset points at this logical level; a boundary there is tested by
(SA3a) on the corresponding segment of the specified controlled circuit,
with the selector included in its register lists. For occurrence $(o,\epsilon)$ of $\mathsf M(R,K)$, let $L_o$ be the
register list at its entry in the expansion of the whole program, let
$U$ and $K_\epsilon$ denote the denotations of $R$ and of the central
operation on $L_o$ (tensor-identity lifts of
Definition~\ref{def:tiera} included), let $F_o\subseteq L_o$, and let
$J_o$ be the ordered product of all expansion operations before its
entry, so that the exact entry subspace is
$\mathcal S_o=J_o\mathcal S_0$. Equation (SA3), with $K_\epsilon$ in
place of $K$, is jointly necessary and sufficient for zero at all
those exits for every allowed initial input. For the keeps of
Definition~\ref{def:keep}, $K^\dagger=K$, but the plus and minus
occurrences can have different entry subspaces. For any specified retained circuit with a specified initial subspace
$\mathcal S_0^{\mathrm{circ}}$ of its own input list (for a clean
implementation of a term, $\mathcal S_0\otimes\ket0_Z$ with workspace
named), the same conjunction assertion holds for a finite collection of
segment exits tested by (SA3a), with each segment's input subspace the
exact image of $\mathcal S_0^{\mathrm{circ}}$ under the prefix ending at
that segment's entry. These assertions include inputs supported on
$\mathcal S_0\otimes\cH_R$ and untouched references.
Replacing exact entry subspaces by containing subspaces preserves
sufficiency; necessity for actual reachable inputs is then not asserted.
No equivalence between final cleanliness and cleanliness of every
syntactic subterm is claimed. Blank-record use requires
zero at its actual entry, supplied by fresh zero allocation, an earlier
proved reset and preservation until use, or another proof of that entry
condition. Neither (SA4) nor a reset alone permits deletion of a factor
still required by an enclosing inverse.
\end{theorem}
\begin{proof}
 \step{sA-rec-reset}{Reset on $\mathcal S$ is equivalent to (SA3).}
 \begin{proof}
  \step{sA-rec-support}{Reset on every density operator supported on
  $\mathcal S$ is equivalent to $(\id-P_F)WP_{\mathcal S}=0$.}
  \begin{proof}\pf\ Test rank-one inputs $\proj\psi$ with
  $\psi\in\mathcal S$. Their output has zero $F$ marginal exactly
  when $W\psi\in\operatorname{ran}P_F$, by the positive-support
  calculation in Theorem~\ref{thm:sA-clean}. This gives the operator
  equality; its converse for mixed inputs follows by a spectral
  decomposition supported on $\mathcal S$.
  \end{proof}
  \qedstep
  \begin{proof}\pf\ Multiply the equality of
  \stepref{sA-rec-support} by $U$ on the left, use
  $U(\id-P_F)=(\id-UP_FU^\dagger)U$ and $UW=KU$.
  Multiplication by $U$ is reversible, proving the equivalence.
  \end{proof}
 \end{proof}
 \step{sA-rec-full}{For $\mathcal S=\operatorname{ran}P_F$ the
 commutator criterion is equivalent.}
 \begin{proof}\pf\ Set $\Pi_F=UP_FU^\dagger$. Equation (SA3)
 becomes $(\id-\Pi_F)K\Pi_F=0$ after right multiplication by
 $U^\dagger$. The finite-dimensional invariant-subspace argument in
 Theorem~\ref{thm:sA-clean} gives $[K,\Pi_F]=0$ and its converse.
 \end{proof}
 \step{sA-rec-marginal}{Marginal restoration is equivalent to (SA4).}
 \begin{proof}
  \step{sA-rec-marginal-need}{\assume{Restoration for all states
  supported on $\mathcal S$.}\prove{Equation (SA4).}}
  \begin{proof}\pf\ Write
  $D=W^\dagger(\id\otimes B_{ij})W-\id\otimes B_{ij}$.
  Equality of the marginals implies $\langle\psi,D\psi\rangle=0$
  for every $\psi\in\mathcal S$, by homogeneity from normalized
  inputs. For $u,v\in\mathcal S$, apply this to $u+v$ and $u+iv$,
  and subtract the diagonal terms: one obtains respectively
  $\langle u,Dv\rangle+\langle v,Du\rangle=0$ and
  $i\langle u,Dv\rangle-i\langle v,Du\rangle=0$.
  Hence both cross terms vanish, without requiring $D$ Hermitian.
  Thus $P_{\mathcal S}DP_{\mathcal S}=0$.
  \end{proof}
  \step{sA-rec-marginal-suff}{\assume{Equation (SA4).}
  \prove{Restoration including every external reference $R$.}}
  \begin{proof}\pf\ For any matrix $C$ on $R$, tensor the compressed
  equality with $C$. Every joint density operator $\sigma$ with support
  in $\mathcal S\otimes\cH_R$ satisfies
  \[
   \tr\!\left[\sigma\bigl(W^\dagger(\id\otimes B_{ij})W
          -\id\otimes B_{ij}\bigr)\otimes C\right]=0.
  \]
  Matrix units $B_{ij}$ and matrix units $C$ span all operators on
  $M'R$. Equality of their expectations is equality of the reduced
  density operators. This proves the claimed restoration.
  \end{proof}
  \qedstep
  \begin{proof}\pf\ Necessity holds by taking the reference
  one-dimensional in \stepref{sA-rec-marginal-need}; sufficiency is
  \stepref{sA-rec-marginal-suff}.
  \end{proof}
 \end{proof}
 \step{sA-rec-segment}{The general segment criterion is (SA3a), also
 with untouched references.}
 \begin{proof}\pf\ For a unit vector $\psi\in\mathcal S$, the
 isometric output $V\psi$ has zero $F$ marginal exactly when
 $(\id_B-P_F^B)V\psi=0$, by the positive-support argument of
 Theorem~\ref{thm:sA-clean}. Testing all such vectors gives (SA3a);
 spectral decomposition proves its converse for density operators.
 Tensoring the zero operator with $\id_R$ gives the same assertion
 on $\mathcal S\otimes\cH_R$. Conversely a one-dimensional reference
 recovers (SA3a). No intermediate workspace condition was used.
 \end{proof}
 \step{sA-rec-prefix}{Exact entry subspaces can be determined without
 assuming any intermediate cleanliness.}
 \begin{proof}
  \step{sA-rec-prefix-rules}{For a fixed-space term define
  $\mathcal T_P(\mathcal S)=\mathcal U_P\mathcal S$. Then
  \[
  \begin{aligned}
   \mathcal T_I(\mathcal S)&=\mathcal S,\\
   \mathcal T_a(\mathcal S)&=A_a\mathcal S,\\
   \mathcal T_{P;R}(\mathcal S)&=\mathcal T_R(\mathcal T_P(\mathcal S)),\\
   \mathcal T_{\mathsf M(P,K)}(\mathcal S)
      &=\mathcal U_P^\dagger K\mathcal U_P\mathcal S.
  \end{aligned}
  \]}
  \begin{proof}\pf\ Substitute the denotations in
  Definition~\ref{def:tiera}, also for the derived menu atom case.
  For an outer allocation use $A_B\mathcal S$; for an adjoint atom
  use $A_a^\dagger\mathcal S$. At an annotated occurrence
  $(o,\epsilon)$ of $\mathsf M(P,K)$, the body starts in
  $\mathcal S_o$, its central operation starts in
  $\mathcal U_P\mathcal S_o$, and the inverse body starts in
  $K_\epsilon\mathcal U_P\mathcal S_o$.
  The expansion rules multiply to $\mathcal U_P^\dagger
  K_\epsilon\mathcal U_P$: for the minus occurrence this is
  $(\mathcal U_P^\dagger K\mathcal U_P)^\dagger$, as proved in
  Theorem~\ref{thm:sA-sound}. The reversed sequence rule likewise
  multiplies to the adjoint. Induction on the term therefore validates
  both signed expansions, including every nested occurrence. Prefixes
  inside an inverse body are evaluated from its actual entry subspace.
  \end{proof}
  \step{sA-rec-prefix-projector}{If a prefix isometry is $V$ and the
  initial subspace is $\mathcal S_0$, its exact entry projector is
  $VP_{\mathcal S_0}V^\dagger$.}
  \begin{proof}\pf\ This operator is self-adjoint and its square is
  itself because $V^\dagger V=\id$. Its range is precisely
  $V\mathcal S_0$. Each logical prefix is a product of total unitary
  atoms, keeps, their adjoints, and zero embeddings; a physical prefix
  is a product of total gates and zero embeddings. Each is an isometry
  by the same multiplication as in Theorem~\ref{thm:sA-sound}.
  Moreover every density operator $\rho$ supported on $V\mathcal S_0$
  is reached from $V^\dagger\rho V$, which has trace one, is supported
  on $\mathcal S_0$, and is taken back to $\rho$ by $V$.
  The same statements hold after tensoring with $\id_R$.
  \end{proof}
  \qedstep
  \begin{proof}\pf\ The syntax and its expansion are finite, so
  \stepref{sA-rec-prefix-rules} and \stepref{sA-rec-prefix-projector}
  determine every requested entry subspace by finitely many operations
  on the specified finite matrices. For a supplied physical circuit,
  take its actual gate prefixes on all registers, including workspace.
  These are exact algebraic prescriptions; no efficient procedure for
  deciding their zero equalities is claimed.
  \end{proof}
 \end{proof}
 \step{sA-rec-all}{The tests have exactly the asserted logical scope.}
\begin{proof}
 \step{sA-rec-all-mark}{\assume{The specified exit is a marked
 occurrence.}\prove{Its (SA3) test with $K_\epsilon$ is necessary
 and sufficient for zero at that exit for every allowed initial state.}}
 \begin{proof}\pf\ Step~\stepref{sA-rec-prefix} identifies the
 actual operator as $U^\dagger K_\epsilon U$ and makes every state
 on the exact entry subspace reachable. Apply
 \stepref{sA-rec-reset} with $K_\epsilon$; Lemma~\ref{lem:keep}
 gives $K_-=K_+$ for a keep. The resulting test is equivalent to
 zero at this exit for every allowed initial state.\end{proof}
 \step{sA-rec-all-segment}{\assume{The specified exit is another
 segment exit.}\prove{Its (SA3a) test on the exact entry subspace
 is necessary and sufficient for zero there for every allowed initial state.}}
 \begin{proof}\pf\ Step~\stepref{sA-rec-segment} gives exactly
 the zero-exit test on its entry subspace. Reachability from
 \stepref{sA-rec-prefix} makes it necessary and sufficient for
 every allowed initial state.\end{proof}
 \step{sA-rec-all-conjunction}{The finite conjunction requires exactly
 that all specified exits be zero for every allowed initial state,
 also in the reference forms.}
 \begin{proof}\pf\ Apply \stepref{sA-rec-all-mark} and
 \stepref{sA-rec-all-segment} separately to each exit. Universal
 quantification over initial states commutes with finite conjunction.
 Reference forms follow from \stepref{sA-rec-segment} and
 \stepref{sA-rec-prefix}.\end{proof}
 \step{sA-rec-all-containing}{\assume{$\mathcal S\subseteq
 \widehat{\mathcal S}$.}\prove{Vanishing on $\widehat{\mathcal S}$
 implies vanishing on $\mathcal S$.}}
 \begin{proof}\pf\ Every vector in $\mathcal S$ is a vector in
 $\widehat{\mathcal S}$, so restriction of the zero operator to
 $\mathcal S$ is zero.\end{proof}
 \step{sA-rec-all-entry}{Forward-entry blankness is an additional
 support condition, and the tests do not specify missing factors
 in later inverses.}
 \begin{proof}\pf\ Definition~\ref{def:clean} requires incoming
 support on the zero-auxiliary subspace. A prior reset establishes
 that support only if intervening operations preserve it. Every
 tested equality is an equality of retained operators, whose
 domains still contain all retained factors. It therefore supplies
 no action on a register list from which a factor has been removed.
 \end{proof}
 \qedstep
 \begin{proof}\pf\ The preceding steps establish the scope and
 entry requirements without assuming intermediate cleanliness.
 \end{proof}
\end{proof}
 \qedstep
 \begin{proof}\pf\ By \stepref{sA-rec-reset}--\stepref{sA-rec-all}.
 \end{proof}
\end{proof}

\begin{counterexample}[Limits of local reset tests]\label{ex:sA-rec-limits}
The following finite Tier~A examples show that (i) replacing a reached
subspace by the whole zero-auxiliary subspace can destroy necessity,
and the commutator criterion need not be necessary on a proper subspace;
(ii) final cleanliness does not imply cleanliness of every inner mark;
(iii) marginal restoration does not permit discarding correlations; and
(iv) a unitary reduced channel need not reset to its initial machine state.
Without an incoming-zero hypothesis, the full-space commutator alone
need not even imply reset on a specified input subspace.
In (i)--(ii) let $Q=r$, $M=a$, and use
\[
 U=\operatorname{CNOT}_{r\to a},\qquad
 K=\operatorname{CNOT}_{a\to r},\qquad
 P_a=\id_r\otimes\proj0_a.
\]
Here and below $K_{0,t}$ denotes the single keep entry with target $t$
and constant-zero Boolean function on the entire machine list.
In (iii) take $Q=q$, $M'=a$, and $\mathcal S$ spanned by
$\Omega=(\ket{00}+\ket{11})_{qa}/\sqrt2$, with $U=\id$ and
$K=K_{0,q}=\id$, so $W=U^\dagger KU=\id$.
In (iv) take $Q=r$, $M=a$ initially zero, body $U$ as above and constant
keep $K=X_r$.
\end{counterexample}
\begin{proof}
 \step{sA-limits-swap}{$U^\dagger KU=\operatorname{SWAP}_{ra}$.}
 \begin{proof}\pf\ On a basis pair $(y,z)$ in order $r,a$, the
 three CNOTs give successively $(y,z\oplus y)$,
 $(z,z\oplus y)$, and $(z,y)$. These are the SWAP basis values;
 linearity gives the operator equality.
 \end{proof}
 \step{sA-limits-subspace}{Item (i): the reached-subspace test need
 not agree with the full zero-subspace test.}
 \begin{proof}\pf\ For $\mathcal S=\operatorname{span}\{\ket{00}\}$,
 SWAP fixes its vector, so reset holds. But
 $\ket{10}\in\operatorname{ran}P_a$ is taken to $\ket{01}$,
 outside that range. Thus reset fails on the containing space
 $\operatorname{ran}P_a$, and Theorem~\ref{thm:sA-recursive}
 gives $[K,UP_aU^\dagger]\ne0$. This also proves that a containing
 subspace test is not necessary for actual reachable inputs.
 For the same SWAP, the subspace spanned by $\ket{01}$ resets to
 $a=0$, but is not initially blank: reset alone does not establish
 the incoming-zero requirement. Separately, take $U=\id$ and
 $K=K_{0,r}=\id$ with
 that same input subspace. Then $[K,UP_aU^\dagger]=0$ while $a$
 remains one, disproving sufficiency of the commutator without the
 incoming-zero hypothesis.
 \end{proof}
 \step{sA-limits-nesting}{Item (ii): final cleanliness does not imply
 inner cleanliness.}
 \begin{proof}\pf\ Let $R=\mathsf M(U,K)$, identifying the displayed
 Boolean gate $U$ with its atom, and let $K_0=K_{0,r}$.
 The nested term $\mathsf M(R,K_0)$ has denotation
 $\mathcal U_R^\dagger\mathcal U_R=\id$. It is therefore clean on
 every $r$ input with $a=0$. Nevertheless, its inner forward term
 $R$ takes input $\ket{10}$ to $\ket{01}$ by
 \stepref{sA-limits-swap}, so that inner mark is not clean.
 No logical cell is discarded or reused as blank in this example.
 \end{proof}
 \step{sA-limits-reference}{Item (iii): marginal restoration does not
 permit discarding correlations.}
 \begin{proof}\pf\ Identity satisfies (SA4) on every subspace and
 leaves the $a$ marginal of $\proj\Omega$ equal to $\id_a/2$.
 Dropping $a$ and allocating a fresh wire $a'$ in zero gives
 $(\id_q/2)\otimes\proj0_{a'}$. After identifying its two-dimensional
 space with that of $a$ solely to compare states, this differs from
 $\proj\Omega$:
 its expectation of $\proj\Omega$ is $1/4$, whereas the original
 expectation is $1$. Even replacement by a state with the same
 marginal $\id_a/2$ gives $\id_{qa}/4$ and the same discrepancy.
 Thus preservation of a marginal does not preserve the joint state
 under deletion and replacement.
 \end{proof}
 \step{sA-limits-unitary}{Item (iv): a unitary reduced channel need
 not reset the machine.}
 \begin{proof}\pf\ With the constant keep $X_r$, conjugation by
 $\operatorname{CNOT}_{r\to a}$ gives $W=X_rX_a$: on a basis pair
 $(y,z)$ the result is $(y\oplus1,z\oplus1)$. Hence
 $W(\psi_r\ket0_a)=X\psi_r\ket1_a$. The reduced channel is
 conjugation by $X$, but $a$ is never returned to zero.
 \end{proof}
 \qedstep
 \begin{proof}\pf\ Steps \stepref{sA-limits-subspace}--\stepref{sA-limits-unitary}
 establish the four assertions. All keeps are of
 Definition~\ref{def:keep}, including constant-zero and constant-one
 Boolean functions, and each body is identity or an explicitly
 displayed Boolean-gate term admitted by Definition~\ref{def:tiera}.
 \end{proof}
\end{proof}

\subsection{A classical sufficient condition and its limits}

The compute--copy--uncompute organization is stated in Bennett's abstract
\cite[abstract, journal p.~525, PDF p.~1]{Bennett1973}.
The following quantum extension and its limitations are proved directly.

\begin{theorem}[Compute--copy--uncompute with an untouched target]
\label{thm:sA-classical}
Let $X,r,M$ be finite pairwise disjoint registers, and let $m_0$ be a
computational-basis configuration of $M$. Suppose $U=U_{XM}\otimes\id_r$
is a total unitary, and for every computational-basis value $x$ of $X$
there are a unit vector $\phi_x\in\cH_X$ and a computational-basis
configuration $m_x$ of $M$ such that, in order $X,M$,
\[
 U_{XM}(\ket x\ket{m_0})=\ket{\phi_x}\ket{m_x}.          \tag{SA5}
\]
Let $v:\{0,1\}^{|M|}\to\{0,1\}^{|r|}$ be a total Boolean function
with the specified finite circuit of Definition~\ref{def:keep}, and
let $K=K_{v,r}$, with $M$ the whole keep-control list. Any additional
machine controls must be included in $M$ and in (SA5).
Then $U^\dagger KU$ is clean and induces the unitary
\[
 \ket{x,y}\longmapsto\ket{x,y\oplus v(m_x)}             \tag{SA6}
\]
on $Xr$, for every quantum input on $Xr$ and any external reference.
The action on inputs with $M$ in $\ket{m_0}$, and hence (SA6),
depends only on (SA5) and independence of $r$, not on the choice of
total extension of $U_{XM}$ away from these inputs. No independence of
the total operator $U^\dagger KU$ on other input states is asserted.

In particular, (SA5) holds with $\phi_x=\ket x$ for a classical
computation preserving the input register $X$. It also
holds whenever $U_{XM}$ is a total computational-basis permutation,
including a fixed finite computation composed only of reversible classical
steps, independent of the unread keep target $r$. Literal preservation
of $X$ is therefore unnecessary for the stronger sufficient condition.

Two enlargements of this universal claim are false. First, allowing
non-basis machine output in (SA5) fails for $X$ empty, $M=a$, $m_0=0$,
$U=H_a\otimes\id_r$, $K=\operatorname{CNOT}_{a\to r}$:
\[
 W\ket{0,0}_{ar}=
   (\ket0_a\ket+_r+\ket1_a\ket-_r)/\sqrt2,
 \qquad\ket\pm=(\ket0\pm\ket1)/\sqrt2.                \tag{SA7}
\]
Second, allowing $U\ket{x,y,m_0}=\ket{x,y,m_{x,y}}$ to depend on the
keep target fails for $M=a$, $U=\operatorname{CNOT}_{r\to a}$,
$K=\operatorname{CNOT}_{a\to r}$: input $\ket{1,0}_{ra}$ ends at
$\ket{0,1}_{ra}$. These are counterexamples to enlarged universal
claims, not necessity claims for each individual clean operation.
For an entangled-output failure upon allowing unrestricted outputs in
place of (SA5), see the separately proved Theorem~\ref{thm:sA-four}.
\end{theorem}
\begin{proof}
 \step{sA-class-basis}{For every basis pair $(x,y)$ the three stages
 have the asserted action.}
 \begin{proof}
  \step{sA-class-compute}{The forward stage gives
  $\ket{\phi_x}_X\ket y_r\ket{m_x}_M$.}
  \begin{proof}\pf\ This is (SA5) and $U$ acting as identity on $r$.
  \end{proof}
  \step{sA-class-copy}{The keep gives
  $\ket{\phi_x}_X\ket{y\oplus v(m_x)}_r\ket{m_x}_M$.}
  \begin{proof}\pf\ The machine factor is one basis configuration,
  so in the sector sum defining $K$ only its $m_x$ summand acts.
  \end{proof}
  \step{sA-class-uncompute}{The inverse gives
  $\ket x_X\ket{y\oplus v(m_x)}_r\ket{m_0}_M$.}
  \begin{proof}\pf\ Total unitarity and (SA5) imply
  $U_{XM}^\dagger(\ket{\phi_x}\ket{m_x})=\ket x\ket{m_0}$.
  Independence of $r$ permits applying this equality to the modified
  target value. No unspecified forward value of $U_{XM}$ is used.
  \end{proof}
  \qedstep
  \begin{proof}\pf\ Compose \stepref{sA-class-compute}--\stepref{sA-class-uncompute}.
  \end{proof}
 \end{proof}
 \step{sA-class-unitary}{The induced operation on $Xr$, with $M$
 initially in $\ket{m_0}$, is unitary on every $Xr$ input and is
 independent of the total extension.}
 \begin{proof}\pf\ For a fixed $x$, XOR by $v(m_x)$ permutes the
 basis of $r$ and squares to identity. Consequently (SA6) is a
 permutation of the full basis of $Xr$. The equality in
 \stepref{sA-class-basis} extends by linearity, also after tensoring
 an identity on a reference. Its expression uses only $m_x$; the
 inverse equality there holds for every unitary extension agreeing
 with (SA5). Thus $M$ is returned to $m_0$ with the claimed induced
 unitary.
 \end{proof}
 \step{sA-class-permutation}{Reversible classical steps with an
 untouched target satisfy the sufficient condition.}
 \begin{proof}\pf\ A finite product of computational-basis
 permutations is a computational-basis permutation. For each input
 $(x,m_0)$ its output is one basis pair $(z_x,m_x)$, so take
 $\phi_x=\ket{z_x}$. This verifies (SA5), whether or not $z_x=x$.
 \end{proof}
 \step{sA-class-H}{Equation (SA7) disproves deletion of the basis-output
 requirement.}
 \begin{proof}\pf\ On $\ket{00}_{ar}$ the forward $H_a$ gives
 $(\ket{00}+\ket{10})/\sqrt2$. The keep gives
 $(\ket{00}+\ket{11})/\sqrt2$. Applying $H_a$ again gives
 $(\ket{00}+\ket{10}+\ket{01}-\ket{11})/2$, which is (SA7).
 Its $a=1$ probability is $1/2$. The forward machine state was the
 non-basis vector $\ket+$, while the keep target was untouched by $U$.
 \end{proof}
 \step{sA-class-target}{The target-dependent example disproves deletion
 of the unread-target requirement.}
 \begin{proof}\pf\ The three CNOTs give
 $\ket{10}_{ra}\mapsto\ket{11}_{ra}\mapsto\ket{01}_{ra}
 \mapsto\ket{01}_{ra}$. Thus $a$ is one at the end although it
 was initially zero. The forward rule is nevertheless the
 basis-valued rule $\ket{y,0}\mapsto\ket{y,y}$.
 \end{proof}
 \qedstep
 \begin{proof}\pf\ By \stepref{sA-class-basis}--\stepref{sA-class-target}.
 The failures arise after deleting the indicated hypothesis; they do
 not contradict the sufficient statement with all hypotheses retained.
 \end{proof}
\end{proof}

\subsection{A four-qubit two-copy circuit, its leakage, and its channels}

\begin{theorem}[Four-qubit calculation]\label{thm:sA-four}
Let $q,r,a,b$ be distinct qubits, with $Q=qr$ and $M=ab$ initially
$\ket{00}$. The copying operations below are the explicit Boolean-gate
atoms of Definition~\ref{def:tiera}, not the full-snapshot source-read
translation of Definition~\ref{def:joint-tick}. Set
\[
 U=\operatorname{CNOT}_{q\to b}H_q\operatorname{CNOT}_{q\to a},
 \quad K=\operatorname{CNOT}_{b\to r},\quad W=U^\dagger KU,
 \quad P_\pm=\proj\pm_r.
\]
Define $E_{00}\xi=\xi_{qr}\otimes\ket{00}_{ab}$. On every input
$\xi\in\cH_{qr}$,
\[
 WE_{00}\xi=B_0\xi\otimes\ket{00}_{ab}
             +B_1\xi\otimes\ket{10}_{ab},\qquad
 B_0=\id_q\otimes P_+,\quad B_1=X_q\otimes P_-.         \tag{SA8}
\]
For any finite untouched reference $R$, the same equality holds for
$\xi\in\cH_{qrR}$ with $W,E_{00},B_0,B_1$ each tensored with
$\id_R$, and the factors canonically reordered by their names.
Consequently the following assertions hold.
\begin{enumerate}
 \item For $r$ initially $0$ and every unit vector $\psi_q$, the final
 state, in order $q,r,a,b$, is
 \[
  \frac{\ket\psi_q\ket+_r\ket0_a+
              X\ket\psi_q\ket-_r\ket1_a}{\sqrt2}
          \otimes\ket0_b.                              \tag{SA9}
 \]
 The probability that $M\ne00$ is $1/2$, independently of the input
 density operator of $q$ and of any reference: the reference input
 is $\rho_{qR}\otimes\proj0_r$, with $\rho_{qR}$ arbitrary.
 The factor $b$ is always zero.
 \item Dropping $M$ gives the channel on $qr$
 \[
  \Phi(\sigma)=B_0\sigma B_0^\dagger+B_1\sigma B_1^\dagger.
                                                               \tag{SA10}
 \]
 It is not a unitary channel on the full space $\cH_{qr}$. For $r=0$,
 \[
  \Phi(\rho\otimes\proj0_r)
   =\tfrac12\rho\otimes P_++\tfrac12X\rho X\otimes P_-.
                                                               \tag{SA11}
 \]
 Every pure $q$ input yields a rank-two $qr$ output of purity $1/2$.
 \item For any fixed density operator $\tau$ initializing $r$, with
 product input $\rho_q\otimes\tau_r$, put
 $p_\pm=\tr(P_\pm\tau)$. Then
 \[
 \begin{split}
  \Phi(\rho\otimes\tau)&=p_+\rho\otimes P_+
                         +p_-X\rho X\otimes P_-,\\
  \tr_r\Phi(\rho\otimes\tau)&=p_+\rho+p_-X\rho X.
 \end{split}                                                  \tag{SA12}
 \]
 The channel in the second line is unitary on $q$ if and only if
 $\tau=P_+$ or $\tau=P_-$. The map from $q$ to retained $qr$ in
 the first line is isometric if and only if the same condition holds;
 the word unitary here applies only after identifying the fixed $r$
 factor with a one-dimensional space. For $\tau=P_+$ the action is
 identity on $q$ and $M$ resets. For $\tau=P_-$ the action is $X$
 on $q$ and $M$ is the fixed state $\ket{10}$, so it does not reset.
 \item When $r=0$, the $q$ marginal is $(\rho+X\rho X)/2$.
 It equals $\rho$ precisely when $X\rho X=\rho$; among pure states
 these are exactly the $X$ eigenstates. This does not make (SA11)
 pure. If the intervening $H_q$ is omitted, the marked operation on
 the same initialized machine is instead CNOT from $q$ to $r$,
 with $M$ reset on every $qr$ input.
\end{enumerate}
Moreover
$U(\ket0_q\ket{00}_{ab})
=(\ket0_q\ket{00}_{ab}+\ket1_q\ket{01}_{ab})/\sqrt2$, which is not a
product vector across $q$ and $ab$; hence, with $X=q$, $M=ab$ and target
$r$ untouched by $U$, unrestricted forward outputs cannot replace (SA5)
in Theorem~\ref{thm:sA-classical}.
The example therefore refutes universal cleanliness for Tier~A, even
with an unread keep target, while preserving its total unitary soundness.
\end{theorem}
\begin{proof}
 \step{sA-four-conjugate}{On the full four-qubit space,
 $W=P_+\otimes\id_{qab}+P_-\otimes X_qX_aZ_b$, with the
 tensor factors placed on their indicated wires.}
 \begin{proof}
  \step{sA-four-K}{$K=P_+\otimes\id_{qab}+P_-\otimes Z_b$,
  with identities on $q,a$ in the second summand.}
  \begin{proof}\pf\ The target Pauli $X_r$ acts by $+1$ on
  $\ket+$ and by $-1$ on $\ket-$. In the latter sector,
  controlled-$X_r$ contributes $(-1)^b$, the action of $Z_b$.
  \end{proof}
  \step{sA-four-Z}{$U^\dagger Z_bU=X_qX_aZ_b$.}
  \begin{proof}\pf\ Conjugating $Z_b$ by
  $\operatorname{CNOT}_{q\to b}$ gives $Z_qZ_b$: on a basis
  pair the multiplier is $(-1)^{b\oplus q}=(-1)^b(-1)^q$.
  Next $H_qZ_qH_q=X_q$ by the matrices in
  Definition~\ref{def:circuit}. Finally conjugating $X_q$ by
  $\operatorname{CNOT}_{q\to a}$ gives $X_qX_a$, because
  the basis pair $(q,a)$ becomes $(q\oplus1,a\oplus1)$.
  The other factors commute with gates on disjoint wires.
  \end{proof}
  \qedstep
  \begin{proof}\pf\ $U$ does not act on $r$, so conjugate
  \stepref{sA-four-K} by $U$ and use \stepref{sA-four-Z}.
  \end{proof}
 \end{proof}
 \step{sA-four-isometry}{Equation (SA8) holds.}
 \begin{proof}\pf\ In \stepref{sA-four-conjugate}, the first
 summand leaves $\ket{00}_{ab}$ fixed. The second takes it to
 $\ket{10}_{ab}$, since $Z_b\ket0_b=\ket0_b$ and
 $X_a\ket0_a=\ket1_a$. The remaining operators are $B_0,B_1$.
 Each identity remains valid after tensoring with a reference identity.
 \end{proof}
 \step{sA-four-state}{Equation (SA9) and the leakage probability hold.}
 \begin{proof}\pf\ Substitute $\xi=\psi\otimes\ket0_r$ in
 (SA8) and use $P_\pm\ket0=\ket\pm/\sqrt2$.
 The two machine vectors in (SA8) are orthogonal, and
 $B_1^\dagger B_1=\id_q\otimes P_-$. Thus the $a=1$
 probability for any $qr$ input $\sigma$ is
 $\tr[(\id_q\otimes P_-)\sigma]$. For $r=0$ this is $1/2$,
 even if $q$ has a reference, since its input trace is one.
 Both summands have $b=0$.
 \end{proof}
 \step{sA-four-channel}{Equations (SA10)--(SA12) hold and define
 trace-preserving completely positive maps.}
 \begin{proof}\pf\ Tracing the orthogonal machine vectors in
 (SA8) removes their cross terms and yields (SA10).
 Moreover $B_0^\dagger B_0+B_1^\dagger B_1
 =\id_q\otimes(P_++P_-)=\id_{qr}$; the Kraus calculation in
 Theorem~\ref{thm:sA-sound} proves the channel properties.
 For every density operator $\tau$ on $r$,
 $P_\pm\tau P_\pm=\tr(P_\pm\tau)P_\pm$ because these
 projectors have rank one. Substitution gives (SA12); taking
 $\tau=\proj0$ gives $p_+=p_-=1/2$ and (SA11).
 \end{proof}
 \step{sA-four-rank}{The full $qr$ channel is not unitary, and every
 pure input on $q$ with $r=0$ has the asserted rank and purity.}
 \begin{proof}\pf\ In (SA11) for $\rho=\proj\psi$, the unit
 vectors $\psi\otimes\ket+$ and $X\psi\otimes\ket-$ are
 orthogonal because their $r$ factors are orthogonal. The two nonzero
 eigenvalues are $1/2,1/2$, so the rank is two and purity is
 $1/4+1/4=1/2$. Unitary conjugation preserves rank-one projections,
 whereas this channel takes the pure input $\proj\psi\otimes\proj0$
 to rank two. It cannot be unitary on $qr$.
 \end{proof}
 \step{sA-four-unitary-when}{The exact unitary and isometric conditions
 for a fixed target input are those in assertion 3.}
 \begin{proof}
  \step{sA-four-extreme}{\case{$p_+=1$ or $p_-=1$.}}
  \begin{proof}\pf\ Positivity of $\tau$ and its zero expectation
 on the complementary projection force its support into the respective
 rank-one eigenspace, by the positive-support argument of
 Theorem~\ref{thm:sA-clean}. Thus $\tau=P_+$ or $P_-$ respectively.
 Equation (SA8) becomes $\psi\ket+\ket{00}$ or
 $X\psi\ket-\ket{10}$ respectively. These are the stated isometries,
 induced unitaries on $q$, and machine outputs.
  \end{proof}
  \step{sA-four-interior}{\case{$0<p_+<1$, so $0<p_-<1$.}}
  \begin{proof}\pf\ For input $\rho=\proj0_q$, the $q$ output
  in (SA12) is $p_+\proj0+p_-\proj1$, of rank two; it cannot
  be a unitary channel. The retained $qr$ output has two nonzero
  eigenvalues $p_+,p_-$ on orthogonal vectors, so the map from $q$
  to $qr$ cannot be isometric either. The cases exhaust all possible
  $p_\pm$ because they are nonnegative and sum to one.
  \end{proof}
  \qedstep
  \begin{proof}\pf\ By \stepref{sA-four-extreme} and
  \stepref{sA-four-interior}.
  \end{proof}
 \end{proof}
 \step{sA-four-fixed}{The fixed states of the $q$ marginal for $r=0$
 are exactly those asserted.}
 \begin{proof}\pf\ The equality $(\rho+X\rho X)/2=\rho$ is
 equivalent to $X\rho X=\rho$. For a rank-one projection this
 says $X\psi=\lambda\psi$ for a scalar of modulus one; applying
 $X$ twice gives $\lambda^2=1$. Conversely either $X$ eigenstate
 satisfies the equality. The rank calculation in
 \stepref{sA-four-rank} still applies to its retained $qr$ output.
 \end{proof}
 \step{sA-four-noH}{Deleting the intermediate Hadamard gives a clean
 CNOT on $qr$.}
 \begin{proof}\pf\ Put $U_0=\operatorname{CNOT}_{q\to b}
 \operatorname{CNOT}_{q\to a}$. On every basis input
 $\ket{x,y,0,0}_{qrab}$, the forward circuit gives
 $\ket{x,y,x,x}$. The keep gives $\ket{x,y\oplus x,x,x}$,
 and the inverse gives $\ket{x,y\oplus x,0,0}$.
 Linear extension proves the claimed action on arbitrary $qr$ inputs.
 \end{proof}
 \step{sA-four-product-limit}{The forward body can entangle data
 and machine while leaving the keep target untouched.}
 \begin{proof}\pf\ On $\ket0_q\ket{00}_{ab}$ its action is
 $(\ket0_q\ket{00}_{ab}+\ket1_q\ket{01}_{ab})/\sqrt2$.
 Tracing $ab$ gives $\id_q/2$, of rank two, whereas a product pure
 vector has rank-one $q$ marginal. Thus this forward vector is not
 of the form (SA5). The operator $U$ never acts on $r$, yet for
 $r=0$ the leakage is $1/2$ by \stepref{sA-four-state}. Hence
 unrestricted forward outputs cannot replace (SA5) in the universal
 cleanliness assertion. This does not claim (SA5) necessary for each
 individually clean operation.
 \end{proof}
 \qedstep
 \begin{proof}\pf\ By \stepref{sA-four-isometry}--\stepref{sA-four-product-limit}.
 The nonzero leakage in \stepref{sA-four-state} disproves universal
 cleanliness, while all factors are total unitaries covered by
 Theorem~\ref{thm:sA-sound}. Their explicit copying atoms do not
 supply an equality with a different source-read translation.
 \end{proof}
\end{proof}
